% SAFA --- copy-edit review copy. Compare against main.tex, which stays canonical.
% Carries the copy editor's edits that survived per-item review, about 155 changes:
% journal style (acronym case, numerals, en-dash percent ranges, closed prefixes),
% factual fixes (table units, leading zeros in Gini, Fig. consistency, 'M s'),
% connectives (so -> thus/hence/therefore, only where the link is a deduction),
% and vocabulary where the meaning is unchanged.
% Rejected: edits that alter claims, delete scope limits, or introduce errors.
% Merge into main.tex once reviewed, then delete this file.
\documentclass{ieeetj}

\usepackage{graphicx}
\usepackage{kotex}
\usepackage{amsmath,amssymb,amsfonts}
\newcommand{\cmark}{\ensuremath{\checkmark}}   % table: public
\newcommand{\xmark}{\ensuremath{\times}}        % table: closed/none
\usepackage{amsthm}
\usepackage{mathrsfs}
\usepackage{booktabs}
\usepackage{multirow}
\usepackage{tabularx}         % Table 1: auto-width wrapping columns (keeps it inside \textwidth)
\newcolumntype{Y}{>{\raggedright\arraybackslash}X}
\usepackage{xcolor}
\usepackage{algorithm}
\usepackage{algpseudocode}
\usepackage{algcompatible}   % uppercase macros \STATE \IF \FOR etc. (Algorithms 1-2)
\renewcommand{\algorithmiccomment}[1]{// #1}  % 교수 코멘트: {..} 대신 // 주석 표기
\providecommand{\RETURN}{\State \Return}  % some TeXLive algcompatible lacks \RETURN -> supply (ignored if defined)
\usepackage{cite}
\usepackage{placeins}         % \FloatBarrier — keep figures inside section/bibliography boundaries
\usepackage{subcaption}       % 2x2 concept diagram (subfigures)
\usepackage[hidelinks]{hyperref}
\usepackage{orcidlink}       % author ORCID icons (\orcidlink)
\AtBeginDocument{\definecolor{tmlcncolor}{cmyk}{0.93,0.59,0.15,0.02}\definecolor{NavyBlue}{RGB}{0,86,125}}
\def\OJlogo{\vspace{-4pt}}
\def\seclogo{\vspace{10pt}}
\def\authorrefmark#1{\ensuremath{^{\textbf{#1}}}}

\newtheorem{theorem}{정리}
\newtheorem{proposition}[theorem]{명제}
\newtheorem{definition}{정의}
\newtheorem{remark}{비고}

% Float spacing tightened to reclaim page space (IEEE defaults are generous).
\setlength{\textfloatsep}{8pt plus 2pt minus 2pt}
\setlength{\dbltextfloatsep}{8pt plus 2pt minus 2pt}
\setlength{\floatsep}{6pt plus 2pt minus 2pt}
\setlength{\dblfloatsep}{6pt plus 2pt minus 2pt}
\setlength{\intextsep}{6pt plus 2pt minus 2pt}
\setlength{\abovecaptionskip}{3pt}
\setlength{\belowcaptionskip}{2pt}

\begin{document}

\receiveddate{XXXX년 XX월}
\reviseddate{XXXX년 XX월}
\accepteddate{XXXX년 XX월}
\publisheddate{XXXX년 XX월}
\currentdate{XXXX년 XX월}
\doiinfo{XXXX.XXXX.0000000}

\markboth{SAFA: 예측 없는 재정렬을 통한 GPU 클러스터 작업 큐의 선두 차단 해소}{Cho 외}

\title{SAFA: 예측 없는 재정렬을 통한 GPU 클러스터 작업 큐의 선두 차단 해소}

\author{Kyunam~Cho\authorrefmark{1,2}\,\orcidlink{0000-0002-8300-8192}, YoungHwan~Jin\authorrefmark{2}\,\orcidlink{0000-0001-9302-9897}, Heonchang~Yu\authorrefmark{1}\,\orcidlink{0000-0003-2216-595X}}
\affil{고려대학교 컴퓨터학과, 서울, 대한민국}
\affil{LG유플러스, 서울, 대한민국}
\corresp{교신저자: Heonchang Yu (e-mail: yuhc@korea.ac.kr).}
\authornote{K. Cho e-mail: mystous@korea.ac.kr. Y. Jin e-mail: yhjin@lguplus.co.kr.}

\begin{abstract}
생산 환경의 GPU 클러스터에서 대부분의 학습 작업은 이제 한 번에 여러 GPU를, 흔히 노드 전체를 요구한다. 이러한 갱 스케줄링 작업은 요청한 모든 GPU가 동시에 비어야만 시작할 수 있다. 크기가 서로 다른 작업이 들어오고 끝나면서 유휴 GPU는 노드들에 흩어지고, 클러스터 전체에는 충분한 유휴 용량이 남아 있어도 큰 작업 하나를 담을 수 있는 단일 노드가 없을 수 있다. 그러면 그 작업은 큐의 선두에서 멈추고 뒤의 모든 작업까지 기다리게 되는데, 이것이 선두 차단(head-of-line, HOL blocking)이다. 그 결과 뒤의 대기 작업이 사용할 수 있는 GPU가 놀게 되고, 큰 작업은 무기한 추월당할 수 있다. 스케줄러의 결정은 두 단계, 즉 다음 작업을 고르는 순서 결정(Order)과 실행 위치를 고르는 배치(Placement)로 나뉜다. 차단된 작업이 뒤의 작업까지 막는지는 순서 결정만이 좌우하므로 HOL 차단은 순서 결정 문제이다. 본 논문은 순서 결정 단계만 변경하는 SAFA(starvation-aware fair aging)를 제안한다. SAFA는 스케줄링 시점에 직접 읽을 수 있는 네 가지 값, 즉 큐에서의 작업 위치, 요청 GPU 수, 이웃 작업과 비교한 대기 연령, 서버별 유휴 GPU 수만을 관찰한다. 아무것도 예측하지 않고, 프로파일링하지 않으며, 최적화 문제를 풀지도 않는다. 대기 작업을 얼마나 앞당길 수 있는지를 계수 하나가 제어하며, SAFA는 그 값을 큐 자체의 통계에서 계산하므로 워크로드가 바뀌어도 다시 튜닝할 필요가 없다. 11만 1천 개 작업으로 구성된 Philly 트레이스를 256~1024 GPU 클러스터에서 재생하였다. 256 GPU에서 여섯 개 순서 결정 정책 모두 순서 공정성 점수가 100점 만점에 0으로 떨어졌으며, 이 중 네 정책은 SAFA가 전혀 받지 않는 실행시간 또는 프로파일 정보를 제공받고도 같은 결과를 보였다. SAFA는 256 및 512 GPU에서 92.7~93.5점을 유지했다. 또한 FIFO 대비 중앙 대기시간을 7\%--13\% 단축하고 평균 GPU 활용률을 92\%에서 95\% 또는 96\%로 높여, 공정성을 처리량 손실 없이 확보했다. 서로 다른 클러스터에서 수집된 Helios와 Alibaba 트레이스에서도 같은 패턴이 나타났다. 실제 NVIDIA B200$\times$8 Kubernetes 클러스터에서는 SAFA를 실제 admission 경로에서 실행했고, 순서 공정성이 차단형 FIFO와 유사한 수준을 보였다. SAFA는 배치 단계에는 전혀 손대지 않으므로 어떤 배치 정책 앞에도 적용할 수 있다. 일곱 개 배치 정책에 대해 실험한 결과, SAFA는 각 정책에서 FIFO 대비 할당률을 3.7~5.8 퍼센트포인트 높였다. 유휴 GPU를 가장 넓게 흩뜨리는 KAI spread에서 향상 폭이 가장 컸다.
\end{abstract}

\begin{IEEEkeywords}
데이터센터 GPU, GPU 갱 스케줄링, 선두 차단, 순서 공정성, 배치 전 재정렬 계층, 예측 없는 스케줄링, 큐 재정렬, 튜닝 없는 스케줄링.
\end{IEEEkeywords}

\maketitle


\section{서론}\label{sec1}
GPU 클러스터의 작업 스케줄링은 운영 비용과 연구 비용에 직접 영향을 미치기 때문에 오랫동안 핵심 문제로 연구되어 왔다. 대규모 조직은 고대역폭 링크로 수백~수천 개 GPU를 연결한 ``하이퍼클러스터''를 구축한다. 서로 다른 GPU 요구량과 작업 완료 시간(job completion time, JCT)을 가진 작업이 연속적으로 들어오면서 클러스터 자원은 단편화된다\cite{athlur2022varuna}. 분할 GPU를 사용하는 공유 환경에서는 이 문제가 더 심각해, 클러스터 전체에 수백 개 GPU가 유휴 상태로 남아 있으면서도 새 작업에는 사용할 수 없는 경우가 발생한다\cite{weng2023beware}. 스케줄링만으로 이 문제를 완전히 해결할 수는 없다. 모든 실행시간을 미리 알고 있어도 대기 작업을 머신에 할당하는 문제는 NP-hard이며\cite{ullman1975np,b1}, 데이터센터 스케줄러는 제출 시점에 실행시간조차 알 수 없다\cite{b2}.

그럼에도 딥러닝 작업 스케줄러는 크게 발전해 왔다. Ye 등\cite{b3}은 관련 연구를 정리하며 현재 최첨단을 대표하는 두 시스템으로, 활용률·공정성·완료 시간을 동시에 개선하는 Sia\cite{b4}와 다중 상태 프로파일링으로 처리량을 높이는 Lucid\cite{b5}를 제시한다. 최근 GPU 스케줄러의 상당수는 이 둘의 변형이며, 의사결정 시점에 직접 관측할 수 없는 값을 예측하거나 프로파일링하거나 최적화 문제로 풀어야 한다. 이러한 정보를 얻으려면 프로파일링 시간, 예측 오차, 전용 인프라 등의 비용이 발생하고, 워크로드가 바뀌면 다시 맞춰야 한다. 각 시스템이 요구하는 정보는 제~\ref{Related}절에서 자세히 설명한다.

이러한 단점을 해결하기 위해 본 연구는 기아 인지 공정 에이징(starvation-aware fair aging, SAFA)을 제안한다. SAFA는 스케줄링 시점에 이미 이용 가능한 값, 즉 대기 작업의 큐 위치, 자원 요구량, 다른 대기 작업과 비교한 상대 연령, 서버별 유휴 GPU 수만으로 결정한다. 여기에는 추정하거나 학습하는 값이 없으므로 실행시간 추정, 처리량 프로파일, 최적화 솔버가 필요 없다. SAFA는 스케줄링 전체를 해결하려 하지 않는다. 자원 단편화가 유발하는 선두 차단(HOL blocking)\cite{karol1987hol}이라는 하나의 실패 모드에 집중한다. HOL 차단은 패킷 스위칭에서 처음 정식으로 다뤄졌으며, 큐 선두의 패킷이 막히면 뒤 패킷까지 모두 정지하는 현상을 뜻한다. GPU 서버 풀은 이 패턴을 그대로 재현한다. 단편화 때문에 큐 선두 작업의 갱을 담을 만큼 큰 서버가 없고, 도착 순서대로만 처리하는 큐는 그 실패를 뒤의 모든 작업에 전달한다. SAFA는 핵심 스케줄러를 바꾸지 않은 채 HOL 차단이 만드는 큐 지연과 도착 순서 불공정을 완화한다. 주요 기여는 다음과 같다.

\begin{enumerate}
\item 무기한 배제를 방지하는 예측 없는 재정렬: SAFA는 HOL로 인한 큐 지연을 줄이되, 그 감소에 필요한 만큼만 도착 순서를 훼손한다. 각 작업이 누적한 연령과 현재 유휴 슬롯에 대한 적합도를 함께 반영하므로 큰 작업이 영구적으로 배제되지 않는다. 현재 슬롯에 잘 맞는 작업은 앞당겨지고, 큰 작업은 대기하면서 연령을 누적해 나머지 큐 대비 우선순위가 상승한다.
\item 배치 정책과의 조합 가능성: SAFA는 어떤 작업을 먼저 처리할지를 정하는 순서 결정 단계만 담당하고, 선택된 작업을 어디에 배치할지는 기존 배치 단계에 맡긴다. 따라서 핵심 스케줄러를 교체하지 않고도 다양한 배치 정책 앞에 동일한 형태로 둘 수 있다. 실험한 일곱 개 배치 정책 모두에서 이 특성이 유지됐으며, 유휴 GPU를 가장 많이 남기는 배치 정책에서 이득이 가장 컸다.
\item 튜닝 없는 계수: 계수 $\alpha$는 HOL 해소와 도착 순서 보존 간의 절충을 제어한다. SAFA는 매 스케줄링 패스마다 큐 연령 통계와 자원 적합도로 이 값을 정규화하므로, 운영자가 조정해야 할 튜닝 노브가 사라지고 워크로드가 바뀌어도 새 그리드 탐색이 필요 없다.
\end{enumerate}

이후 구성은 다음과 같다. 제~\ref{Related}절은 관련 연구를 정리한다. 제~\ref{motivation}절은 큐 재정렬이 필요한 이유를 설명한다. 제~\ref{algo:intro}절은 큐 재정렬의 수학적 정식화와 튜닝 없는 균형 메커니즘을 제시한다. 제~\ref{sec:eval}절은 공개 트레이스(Philly, Helios, Alibaba)와 실제 B200 Kubernetes 측정으로 SAFA를 평가한다. 마지막으로 제~\ref{conclu}절에서 결과를 요약하고 향후 연구 방향을 제시한다.

본 논문에서 ``GPU''는 데이터센터 GPU를 뜻하며, 간결성을 위해 이후에는 이 표현을 그대로 사용한다.

\section{관련 연구}\label{Related}
SAFA와 기존 연구의 위치를 비교하기 위해 두 차원을 구분한다. 첫째는 스케줄러가 개입하는 단계다. 스케줄링을 여러 방식으로 분해할 수 있지만, 본 논문에서 중요한 단계는 어떤 작업을 먼저 처리할지를 정하는 순서 결정(Order)과 선택된 작업을 어디에 놓을지를 정하는 배치(Placement)이다. 둘째는 의사결정에 필요한 값이 모두 결정 시점에 관측 가능한지 여부다. 큐 위치, 자원 요구량, 누적 대기, 관측된 누적 사용량, 현재 노드 점유 상태는 모두 즉시 확인 가능하다. 그 밖의 값은 예측·프로파일링·최적화 계산으로 만들어야 하며, 본 논문에서는 이를 세 범주로 나눈다. 즉 실행시간 추정 또는 순위 예측(SJF, EASY backfilling, learning-to-rank), 작업별 성능 또는 처리량 프로파일(Optimus, Lucid, Sia, Pollux, Arena, Rubick), 그리고 라운드 솔버 또는 경매 입찰(Themis, FFT, RLTune)이다.

본 논문에서 반복적으로 등장하는 두 종류의 FIFO를 구분해야 한다. 차단형 FIFO(blocking FIFO)는 작업을 도착 순서대로 엄격히 처리하므로 선두 작업을 배치할 수 없으면 큐 전체가 멈춘다. 비차단형 FIFO(nonblocking FIFO)는 막힌 선두를 건너뛰고 배치 가능한 첫 작업을 실행한다. 후자는 전자가 낭비하는 유휴 GPU를 회수하지만 더 이상 도착 순서를 보존하지 않으며, 반복해서 건너뛰는 큰 작업은 무기한 기아 상태에 빠질 수 있다.

\begin{table*}[htb!]
\centering
\caption{스케줄러 분류}
\label{tab:related}
\scriptsize
\setlength{\tabcolsep}{4pt}
\begin{tabularx}{\textwidth}{@{}llYcY@{}}
\toprule
\textbf{시스템} & \textbf{발표처} & \textbf{결정 축(정보 가정)} & \textbf{실험} & \textbf{선정 이유 / 대표 기준선} \\
\midrule
\multicolumn{5}{@{}l}{\textit{실험 기준선 --- 동일 계층, 정보 예산이 SAFA보다 낮지 않음}}\\
FIFO (차단형) & --- & 큐 순서(관측값만) & $\checkmark$ & 제어군(순수 정책 효과)\\
SJF & 고전적 & 큐 순서(oracle 실행시간) & $\checkmark$ & 평균 최적 참고선\\
LAS (Tiresias)\cite{gu2019tiresias} & NSDI'19 & 큐 순서(누적 서비스량) & $\checkmark$ & 최소 누적 서비스량\\
Kueue\cite{kueue2024} & k8s & quota gate(관측값만) & $\checkmark$ & 운영 표준\\
EASY\cite{mualem2001backfilling} & HPC & 예약+backfill(oracle 실행시간) & $\checkmark$ & 완전 추정값 변형(상한선은 아님)\\
Themis\cite{mahajan2020themis} & NSDI'20 & FTF 순서(oracle 실행시간, $\pm20\%$; 솔버 생략) & $\checkmark$(변형) & 공정성 축; 라운드 경매는 재현하지 않음\\
Lucid\cite{b5} & ASPLOS'23 & 프로파일 기반 packing & $\checkmark$(변형) & 합성 활용률, 실행시간 oracle 주입 --- 단일 스칼라로 합성 왜곡을 제한\\
\textbf{SAFA} & --- & 큐 순서 + 노드 점유(관측값만) & $\checkmark$ & 본 연구(예측 없음)\\
\midrule
\multicolumn{5}{@{}l}{\textit{직교하는 배치 축 --- SAFA처럼 예측이 없으며 비교가 아니라 조합 대상으로 사용}}\\
FGD\cite{weng2023beware} & ATC'23 & 배치(단편화 회피; 관측값만) & $\checkmark$ & 상보적 단계; SAFA$\oplus$FGD 조합\\
KAI\cite{kai2025} & k8s & 배치(binpack/spread; 관측값만) & $\checkmark$ & 운영 표준; SAFA$\oplus$KAI 조합\\
\midrule
\multicolumn{5}{@{}l}{\textit{관련 연구 --- 트레이스에 결정 변수가 없거나 계층이 다름}}\\
Learning-to-Rank\cite{fu2024efficient} & NeurIPS'24 & LLM 요청 순서(예측 순위, ms) & --- & 계층 불일치; SJF로 대표\\
NotebookOS\cite{notebookos2026} & ASPLOS'26 & 대화형 복제 & --- & 다른 워크로드(범위 밖)\\
Sia\cite{b4} & SOSP'23 & 탄력적 ILP(goodput 프로파일) & --- & 트레이스에 프로파일 결정 변수 없음\\
Gavel\cite{narayanan2020gavel} & OSDI'20 & 이기종 인지 라운드 정책(처리량 행렬) & --- & 처리량 행렬 결정 변수 없음\\
Optimus\cite{peng2018optimus} & EuroSys'18 & 탄력적 worker/PS(성능 모델) & --- & 성능 모델 결정 변수 없음\\
Pollux\cite{qiao2021pollux} & OSDI'21 & 탄력적 goodput(프로파일) & --- & goodput 프로파일 없음\\
Arena\cite{arena2026} & EuroSys'26 & 적응형 병렬화(프로파일+추정) & --- & 실행 계획 결정 변수 없음\\
Rubick\cite{rubick2025} & MLSys'25 & 실행 계획 재구성 & --- & 실행 계획 결정 변수 없음\\
RLTune\cite{rltune2025} & SoCC'25 & RL 순서+MILP(학습 모델) & --- & 학습+솔버; Lucid/Themis로 대표\\
FFT\cite{fft2025} & ICS'25 & 라운드 비용 최소화(epoch 추정) & --- & 솔버+추정; Themis(FTF)로 대표\\
Tetris\cite{grandl2014multi} & SIGCOMM'14 & 다차원 자원 정렬(순서$+$배치) & --- & 결정 변수 없음; 배치와 결합\\
Maui\cite{jackson2001maui} & JSSPP'01 & 다요소 우선순위(실행시간, fairshare) & --- & 추정값과 사용 이력 필요; 잔여 aging은 FIFO로 환원\\
\bottomrule
\end{tabularx}
\end{table*}

\subsection{스케줄링 단계별 기법}
먼저 순서 결정 단계를 살펴본다. Kueue\cite{kueue2024}는 작업을 일시 중단하고 재개하며 quota 기반 공정 공유를 제공하는 Kubernetes 네이티브 큐잉 시스템이다. 핵심 스케줄러를 교체하지 않는 gating 계층으로 동작하기 때문에 본 논문의 운영 환경 기준선과 가장 가깝다. 하지만 공정성 메커니즘은 quota 분배에 한정되며 HOL 기아나 단편화를 직접 다루지 않는다. 또한 quota 공정성은 여러 Virtual Cluster(VC)가 경쟁할 때만 의미가 있고, 단일 VC 내부에서는 Kueue가 비차단형 FIFO로 환원된다. Tiresias\cite{gu2019tiresias}의 LAS는 누적 자원 사용량이 가장 적은 작업을 우선하는 예측 없는 규율로, 다단계 피드백 큐 및 선점과 결합해 평균 JCT를 줄인다. 그러나 비선점 GPU 운영에서는 대기 작업의 누적 자원 사용량이 모두 0으로 유지되므로 순서 키가 정보를 제공하지 못하며, LAS 역시 비차단형 FIFO와 동일해진다. 이 동작은 선점이 핵심인 원래 Tiresias 설정과 다를 수 있다.

배치 축에는 두 방법이 있다. FGD\cite{weng2023beware}는 부분 GPU 점유로 심하게 단편화된 공유 클러스터를 대상으로 하며(트레이스에서 21\%--42\%), 무작위로 선택한 작업이 스케줄되지 못할 확률을 추정한 뒤 그 확률 증가가 가장 작은 노드에 작업을 배치하고, 미할당 GPU를 최대 49\%까지 줄인다. SAFA와 마찬가지로 큐 상태와 노드 점유만 사용하고 예측하지 않는다. 차이는 개입 단계다. FGD는 작업이 선택된 이후 그 작업을 어디에 둘지를 정하므로, 어떤 노드도 갱 전체를 담을 수 없는 선두 작업에는 도움을 주지 못한다. 바로 이런 경우를 SAFA가 다루므로 둘은 경쟁 관계가 아니라 상보 관계이며, 본 논문에서는 FGD를 조합 대상으로 평가한다. 운영 스케줄러 KAI\cite{kai2025}는 binpack과 spread 모드를 제공한다. 기본 binpack은 유휴 GPU가 적은 노드에 더 높은 점수($\text{score}\propto 1-(\text{free}-\min)/(\max-\min)$)를 주어 작업을 집약하고 큰 갱을 위한 빈 노드를 남긴다.

\subsection{파생 정보에 의존하는 기법}
결정 시점에 직접 알 수 없는 값에 의존하는 기법은 세 범주로 나눌 수 있다. (i) 실행시간 추정과 순위 예측, (ii) 작업별 성능 또는 처리량 프로파일, (iii) 라운드 솔버와 경매 입찰이다. 추정 오차, 프로파일링 오버헤드, 전용 인프라를 비용으로 본다면 뒤로 갈수록 획득 비용이 커진다.

(i) 실행시간 추정과 순위 예측: EASY backfilling\cite{mualem2001backfilling}은 표준 HPC 기법으로, 큐 선두에 자원을 예약한 상태에서 그 예약을 방해하지 않는 후속 작업만 앞당긴다. 이 방법은 실행시간 추정이 필요하며, 사용자가 제출하는 실행시간 추정치는 실제와 2--10$\times$ 차이 나는 경우가 흔하다. Learning-to-Rank\cite{fu2024efficient}은 출력 길이의 상대 순위만으로 LLM 서빙에서 SJF/SRTF를 근사하지만, 아직 관측되지 않은 미래 정보인 출력 길이를 예측해야 하며 요청(반복) 단위에서 동작한다.

(ii) 작업별 성능 또는 처리량 프로파일: Lucid\cite{b5}는 작업별 활용률 프로파일에 따라 작은 작업을 함께 배치해 평균 JCT를 줄인다. Optimus\cite{peng2018optimus}, Gavel\cite{narayanan2020gavel}, Sia\cite{b4}, Pollux\cite{qiao2021pollux}, Arena\cite{arena2026}, Rubick\cite{rubick2025} 같은 탄력적 프로파일 기반 스케줄러는 각 라운드에서 GPU 수 또는 종류, 배치 크기, 병렬화 계획을 최적화하며 작업별 goodput 또는 처리량 프로파일이 필요하다. 이런 프로파일링에는 비용과 오차가 모두 존재한다. 선형 처리량 외삽은 최대 2.12$\times$의 오차를 낼 수 있고, 프로파일을 생략하면 JCT가 13\%--20\% 감소한 사례도 보고되었다\cite{schedmate2025}. NotebookOS\cite{notebookos2026}는 대화형 노트북을 위한 온디맨드 복제를 스케줄링하며, 본 논문의 배치 학습 워크로드와 본질적으로 다르다.

(iii) 라운드 솔버와 경매: Themis\cite{mahajan2020themis}는 장기 공정성 지표로 완료시간 공정성 $r$을 도입한다. 이는 공유 클러스터에서의 작업 완료 시간을 전용 $1/N$ 클러스터에서의 완료 시간으로 나눈 값이다. 중앙 ARBITER가 라운드 경매를 통해 최대 $r$을 최소화하며, 애플리케이션은 추정된 $r$로 입찰하고 전략적 조작 불가능성(strategy-proofness)과 Pareto 효율성을 보장한다. 그러나 이 방법은 솔버 기반 경매(경합 시 p95 지연 1,398\,ms)와 애플리케이션이 제출한 $r$ 추정에 의존한다. FFT\cite{fft2025}와 RLTune\cite{rltune2025}은 각각 라운드 비용 최소화 솔버($+$epoch 추정)와 학습된 RL 우선순위($+$MILP 노드 매핑)를 전제로 한다.

\subsection{기준선 선정 근거}\label{baseline-rationale}
비교가 의미 있으려면 SAFA와 기준선이 (i) 동일한 결정 변수, (ii) 동일한 스케줄링 계층을 공유하고, (iii) 정보 예산이 SAFA에 유리해서는 안 된다. 조건 (iii)은 예측형 정책을 어떻게 다룰지 결정하게 한다. SJF, EASY, Themis, Lucid는 작업 실행시간 또는 프로파일을 필요로 하지만, 이 값은 실제로 작업이 실행되기 전에는 존재하지 않는다. 이들을 제외하면 ``예측 없이도 재정렬할 수 있다''는 본 논문의 주장에 예측형 상대가 사라진다. 따라서 트레이스가 사후 기록한 실행시간을 제공해 누락 정보를 보충한다. 오차 없는 값을 주면 추정기 자체의 영향을 제거할 수 있으므로, 각 정책은 자신에게 가장 유리한 정보 조건에서 정책 본연의 동작을 보이게 된다. SJF와 EASY에는 정확한 기록 실행시간을 제공하고, Themis에는 여기에 20\% 무작위 오차를 추가한다. FIFO, LAS, Kueue는 SAFA와 마찬가지로 스케줄링 시점에 제공되는 정보만으로 결정한다. Lucid는 경계에 있다. 트레이스에 활용률 프로파일이 없지만, 이 트레이스에서 재현 가능한 최신 스케줄러 중 하나이므로 제~\ref{sec:eval}-\ref{eval:impl}절에 설명한 단일 스칼라로 왜곡을 제한한 재구성 프로파일을 사용한다. 표~\ref{tab:related}은 비교 정책이 어떤 정보를 받는지 정리한다.

나머지 스케줄러는 조건 (i) 또는 (ii)를 위반하므로 기준선에서 제외한다. (a) oracle로도 공급할 수 없는 결정 변수: Sia, Gavel, Optimus, Pollux, Arena, Rubick이 요구하는 goodput 프로파일, 성능 모델, 재구성 곡선은 트레이스에서 실제로 실행된 적 없는 실행 구성을 나타내므로 기록값으로 공급할 수 없다. 이를 임의로 합성하면 실험은 해당 스케줄러가 아니라 우리가 만든 스케줄러 모델을 평가하게 된다. Sia를 rigid 설정으로 제한하면 FIFO/SJF와 동등해진다. (b) 계층 불일치: Learning-to-Rank(ms, 요청 단위)과 NotebookOS(대화형 복제)는 결정 대상과 시간 규모 모두 클러스터 수준 갱 작업과 다르다. 전자의 원리는 완전 예측 한계인 SJF로 대표한다. (c) 학습 또는 솔버 의존: RLTune과 FFT는 학습 모델 및 라운드 솔버($+$epoch 추정)를 전제로 하므로 예측하지 않는 SAFA와 다른 정보·인프라 계층에 있다. (d) 순서와 자원 적합도 또는 에이징을 함께 쓰는 고전적 선행 연구: Tetris\cite{grandl2014multi}와 Maui\cite{jackson2001maui}가 SAFA와 가장 가깝지만, Maui의 fairshare는 단일 VC 트레이스에 없는 다중 VC 사용 이력이 필요하다. 이를 빼면 에이징은 대기시간만으로 결정된다. 모든 대기 작업이 같은 속도로 나이를 먹으므로 순위는 도착 순서에서 벗어나지 않고 FIFO가 이미 이를 대표한다. Tetris의 다차원 정렬은 전체 GPU 단위 트레이스가 갖고 있지 않은 결정 변수를 요구한다. SAFA의 차별점, 즉 갱 자원 적합도 $R$과 큐 통계로 정규화된 튜닝 없는 계수는 제~\ref{motivation}절에서 동기를 설명하고 제~\ref{algo:intro}절에서 정식화한다. KAI와 FGD는 배치 단계에서 동작하므로 기준선이 아니라 SAFA의 조합 대상으로 평가한다.

\section{동기}\label{motivation}
딥러닝 학습은 이제 거의 전부 컨테이너 안에서 수행되며\cite{b7}, Docker Hub\cite{b8}와 Hugging Face\cite{b9} 같은 레지스트리 덕분에 사전 구축된 프레임워크, 라이브러리, 아키텍처를 쉽게 공유하고 재현할 수 있다. 하지만 이 편의성은 하드웨어 경계에서 비용을 만든다. 컨테이너는 GPU를 완전한 장치 단위로 점유하므로, 스케줄러가 보는 것은 계산량의 풀이 아니라 분할할 수 없는 카드들의 풀이다.

타임 슬라이싱을 쓰면 할당률을 높일 수 있지만 학습 중 이를 유연하게 활용하기는 여전히 어렵기 때문에 배타적 점유가 일반적이다. 모델 규모가 커지면서 하나의 작업이 여러 GPU를 동시에 요구하는 것도 흔해졌다. 완료 시간과 GPU 요구량이 서로 다른 작업이 지속적으로 도착하는 상황에서 배타적 점유는 곧 단편화로 이어진다.

\begin{figure}[htb!]
\centering
\begin{subfigure}{0.8\columnwidth}\includegraphics[width=\linewidth]{Pic/Fig_01.png}\caption{스케줄링 전 서버 풀(단편화 상태)}\label{fig:concept-a}\end{subfigure}

\medskip
\begin{subfigure}{0.8\columnwidth}\includegraphics[width=\linewidth]{Pic/Fig_02.png}\caption{같은 시점의 대기 큐}\label{fig:concept-b}\end{subfigure}

\medskip
\begin{subfigure}{0.8\columnwidth}\includegraphics[width=\linewidth]{Pic/Fig_03.png}\caption{스케줄링 후, $t_2$가 선두에서 차단된 상태}\label{fig:concept-c}\end{subfigure}

\medskip
\begin{subfigure}{0.8\columnwidth}\includegraphics[width=\linewidth]{Pic/Fig_04.png}\caption{남은 대기 큐($t_2$부터 $t_5$까지의 백로그)}\label{fig:concept-d}\end{subfigure}
\caption{단편화와 HOL 차단}
\label{fig:concept}
\end{figure}

그림~\ref{fig:concept}(a)와 (b)는 각각 GPU 서버 풀과 작업 대기 큐의 상태를 보여준다. 이 상태에서 most-allocated 스케줄러는 $\textit{t}_{0}$와 $\textit{t}_{1}$을 배치할 수 있다. 그러나 $\textit{t}_{2}$는 그림~\ref{fig:concept}(c)처럼 다른 작업이 끝나 GPU 슬롯을 반환할 때까지 배치할 수 없다. 큐가 제~\ref{Related}절의 차단형 FIFO처럼 엄격한 도착 순서로 처리되므로, 뒤의 $\textit{t}_{3}$, $\textit{t}_{4}$, $\textit{t}_{5}$는 $\textit{t}_{2}$가 배치된 뒤에야 스케줄될 수 있다. 그림~\ref{fig:concept}(d)는 남은 큐를 보여준다.

문제는 이 시점에 서버 풀 전체의 유휴 GPU 슬롯 수만 보면 이미 $\textit{t}_{2}$와 $\textit{t}_{3}$의 요구량 합을 만족한다는 점이다. 이것이 바로 이미 배치된 작업이 GPU 서버 풀을 단편화해 만드는 HOL 차단이다. 이러한 단편화는 운영 중 지속적으로 생성되고 해소되며, 최악의 경우는 선두에서 막힌 작업이 서버가 가진 $M$개의 가속기 전부를 요구하는데 모든 서버에 $M-1$개씩만 비어 있는 상황이다.

원인은 할당이 공간적으로 이루어지고 GPU가 분할 불가능하기 때문이다. 또한 학습 작업은 체크포인트 저장·재적재 비용이 높고 자원 풀을 재구성하는 것이 불가능하거나 지나치게 비싸므로, 비선점이 사실상 표준이다. 따라서 기아 상태는 작업을 담을 만큼 연속된 큰 슬롯이 생겨야만 해소된다. 스케줄러가 그 슬롯을 더 빨리 만들 수는 없다. 다만 큰 작업을 큐 앞쪽에 계속 유지해, 슬롯이 열리는 순간 그 작업이 곧바로 사용할 수 있도록 할 수 있다.

HOL 차단에는 효율성 비용이 따른다. 선두가 막힌 동안 뒤의 작업이 사용할 수 있는 GPU가 유휴 상태로 남는다. 수백~수천 GPU 규모의 하이퍼클러스터에서는 활용률 몇 퍼센트포인트 차이도 항상 수십 개 GPU가 놀고 있다는 뜻이다. 막힌 선두를 건너뛰면 이 자원을 회수할 수 있지만, 건너뛴 작업은 더 뒤로 밀리고 특히 큰 작업이라면 영원히 실행되지 않을 수 있다. 즉 낭비를 없애는 방식은 큰 작업을 굶기고, 큰 작업을 보호하기 위해 차단을 허용하면 GPU를 낭비한다. 아직 해결되지 않은 것은 차단 자체보다 이 절충이다.

이 절충은 새로운 발견도, 대응 방식도 아니다. 공간 공유 병렬 작업 스케줄링은 JSSPP 계보\cite{feitelson1996toward} 이래 단편화되는 풀에서 분할 불가능하고 비선점인 작업을 전제로 해 왔고, 그 환경에서 기아의 전통적 해법은 에이징이었다. Maui의 XFactor 및 대기시간 가중치\cite{jackson2001maui}, Slurm의 age factor는 대기 시간이 늘수록 우선순위를 높이고, Tetris\cite{grandl2014multi}는 다차원 자원 정렬 점수까지 같은 결정에 포함한다. 하지만 이들 어느 것도 절충을 제거하지 못한다. 모두 운영자가 조정해야 하는 가중치 노브로 절충을 노출시키고, 적절한 값은 워크로드가 바뀌면 함께 이동한다. 기존 에이징은 또한 나이를 먹은 작업이 현재의 유휴 슬롯에 실제로 들어갈 수 있는지 고려하지 않는다. 그 결과 풀에 아직 담을 수 없는 작업을 승격하는 일이 발생할 수 있다. 고전적 접근은 절충을 해결한 것이 아니라 튜닝 파라미터로 옮겨 놓은 셈이다.

예측도 이를 해결하지 못한다. 도착 시각과 완료 시각은 스케줄링 시점에 알 수 없고, 비예지(nonclairvoyant) 설정에는 자체적인 competitive-ratio 하한이 존재한다\cite{motwani1994nonclairvoyant}. 따라서 남은 문제는 HOL 차단이 발생한 뒤에 개입하고, 스케줄링 순간에 관측 가능한 정보만 읽으며, 에이징 가중치를 운영자 노브가 아니라 큐 통계에서 자체 계산하는 메커니즘이다. 제~\ref{algo:intro}절에서 이를 구성하고 제~\ref{sec:eval}-\ref{eval:robust:placement}절에서 배치 정책이 바뀌어도 이득이 유지되는지 확인한다.

\section{작업 큐의 선두 차단 해소}\label{algo:intro}
본 논문에서 갱(gang)은 요청한 GPU 수를 한 번에 전부 받아야 하는 고정 요청이며, Feitelson과 Rudolph의 분류\cite{feitelson1996toward}에서 말하는 rigid job에 해당한다. 여기서 gang이라는 용어는 시간 슬라이싱을 동기화하는 Ousterhout의 gang scheduling\cite{ousterhout1982gang}을 뜻하지 않는다.

\subsection{SAFA 큐 재정렬}
뒤에 도착한 작업에게 계속 추월당하는 작업은 영원히 차례를 얻지 못한다. 중앙 대기시간은 이 현상을 드러내지 못한다. 중앙값은 대개 잘 들어맞는 많은 작업을 반영하고, 들어맞지 않는 소수 작업을 반영하지 않기 때문이다. 따라서 정책이 낮은 중앙값을 보고하면서도 일부 작업을 영원히 실행하지 못할 수 있다. 순서 공정성은 단순한 부가 기능이 아니라 운영 요구사항이다. 어떤 작업도 무기한 배제되어서는 안 된다.

큐 재정렬의 목적은 처리량 그 자체가 아니다. HOL 차단이 뒤 작업에 부과하는 큐 지연을 줄이되, 그 감소에 필요한 만큼만 도착 순서를 잃는 것이다. 백로그가 커질 때 큐가 막힌 작업 하나 뒤에서 계속 정지해 있는 상태를 방지하는 것이 목표다. SAFA는 특정 큰 작업의 완료를 보장하지 않는다. 여전히 충분히 큰 슬롯이 실제로 열려야 하기 때문이다. 제한 없이 작업을 승격하면 공정성을 없애는 대신 다른 곳으로 옮길 뿐이므로, 하나의 가중치가 누적 연령이 도착 순서를 얼마나 강하게 덮어쓸 수 있는지 제어한다.

SAFA는 매 스케줄링 패스에서 순서대로 실행되는 알고리즘~\ref{alg:priority}와~\ref{alg:reorder}로 대기 큐를 재정렬한다. 각 패스 전 큐는 도착 순서로 복원된다. 알고리즘~\ref{alg:priority}은 두 블록으로 조정 우선순위를 계산한다. 첫 블록은 서버 풀을 순회하면서 각 요청 크기의 자원 적합도 $R$을 갱신한다. 두 번째 블록은 대기 작업의 상대 연령을 다시 기준화하고 각 작업의 조정 우선순위 $P^*$를 계산한다. 알고리즘~\ref{alg:reorder}는 가장 큰 $P^*$를 가진 작업을 큐 선두로 이동한다.

자원 적합도 $R_j$는 $j$-GPU 요청이 현재 서버 풀에 얼마나 잘 들어맞는지를 요약한다. 유휴 GPU 수가 $f_{S_i}$인 각 서버 $S_i$에 대해, $j\in\{1,\dots,M_i\}$이고 요청이 들어맞는 경우($j\le f_{S_i}$) $R_j=1$이다. 들어맞지 않으면 부족한 GPU 하나당 0.1을 차감하고 최솟값은 0으로 자른다. 모든 서버 가운데 가장 높은 값을 유지한다. 적합도는 한 서버가 수용할 수 있는 최대 갱 크기 $M=\max_i M_i$까지 모든 요청 크기에 대해 정의된다. 따라서 $R_j\in[0,1]$이며, 요청량이 $a_j$인 대기 작업 $t_j$는 $R_{a_j}$를 사용한다.

두 번째 블록은 우선순위를 계산한다. 연령 $A_j$는 새 작업이 제출될 때마다 1씩 증가하는 음이 아닌 정수이므로, 경과 시간 대신 도착 횟수를 센다. 큐의 최소 연령을 빼면 상대 연령 $\tilde A_j=A_j-\min_k A_k$가 되고, 각 대기 작업의 $P^*$는 다음과 같다.

\begin{equation}
P^*_j = P_j + \alpha\, \tilde A_j \cdot R_{a_j}
\label{eq:01}
\vspace{2pt}
\end{equation}

첫 항 $P_j=1/\text{base}^{\,pos_j}$는 도착 우선순위다. $\text{base}=2$이면 위치 $pos_j=0,1,2$에서 각각 $1$, $0.5$, $0.25$가 되고 큰 값이 우선한다. 두 번째 항 $\alpha\,\tilde A_j R_{a_j}$는 작업이 얼마나 오래 기다렸는지와 현재 빈 슬롯에 얼마나 잘 들어맞는지를 곱한 우선순위 조정값이다. 계수 $\alpha$는 이 곱이 도착 순서를 얼마나 빠르게 덮어쓸지를 정한다. 이를 수치로 보면, 도착이 충분히 일정해 연령 1단위가 10분에 해당한다고 가정하고 큐의 두 번째 작업($P_j=0.5$에서 시작)을 생각할 수 있다. $R_{a_j}=1$일 때 6시간 안에 0.5를 추가하려면 $0.013889\times36\times1\approx0.5$이므로 $\alpha\ge0.5/36\approx0.013889$가 필요하다.

다만 이 예시는 선두 작업 자체의 조정을 무시한 단순화다. 실제로 선두에도 $\alpha\tilde A_0R_{a_0}$가 더해지므로, 작업 $j$가 선두를 추월하려면 $\alpha(\tilde A_jR_{a_j}-\tilde A_0R_{a_0})>P_0-P_j$를 만족해야 한다.

알고리즘~\ref{alg:reorder}은 가장 큰 $P^*$의 위치 $i_{\max}$를 찾고, 그 작업이 이미 선두가 아니라면 다른 작업의 상대 순서는 그대로 둔 채 해당 작업만 맨 앞으로 옮긴다. 각 패스에서 조정 우선순위가 가장 높은 작업 하나만 선두로 이동시키는 이유는, 그 작업을 실제 배치하면 서버 풀의 유휴 슬롯이 바뀌고 그에 따라 $R$과 나머지 작업의 $P^*$가 모두 무효화되기 때문이다. 선택된 작업은 큐 위치, 상대 연령, 자원 적합도를 함께 고려했을 때 가장 높은 조정 우선순위를 갖는다.

$pos_j$와 $\tilde A_j$는 큐에서 직접 읽고, $R_{a_j}$는 작업의 가속기 요청량과 서버별 유휴 슬롯 수를 하나의 적합도 값으로 결합한다. 모두 스케줄링 시점에 읽으며 추정하지 않는다. 선형식은 형태만 보면 Kleinrock의 시간 의존 우선순위\cite{kleinrock1976queueing}와 유사하지만 가정은 다르다. Kleinrock의 보존 법칙은 선점, 단일 서버, 정상상태를 가정하는 반면 비선점·공간 공유·과부하 갱 환경에는 이 가정이 성립하지 않는다. SAFA가 이 고전적 형태에 추가하는 요소는 분할 불가능한 갱의 자원 적합도 $R$과 큐 통계로 정규화된 계수다.

\begin{algorithm}
\caption{SAFA 조정 우선순위}\label{alg:priority}
\textbf{입력:} \\
\hspace*{4mm} 서버 $S_i$ ($i{<}n$): 가속기 수 $M_i$, 유휴 수 $f_{S_i}$ \\
\hspace*{4mm} 대기 작업 $t_i$ ($i{<}l_{wq}$): 가속기 요청량 $a_i$, 연령 $A_i$ \\
\hspace*{4mm} 자원 적합도 $R_j$ ($j\in\{1,\ldots,M\}$) \\
\hspace*{4mm} 계수 $\alpha$ (고정값 또는 식~\eqref{eq:auto-alpha}의 $\alpha_{\text{eff}}$)
\begin{flushleft}
\textbf{출력:} 조정 우선순위 $P^*$
\end{flushleft}
\begin{algorithmic}[1]
\STATE $R \gets 0$
\STATE {}
\FOR{ $i \gets 0$ \textbf{부터} $n-1$}
    \FOR{$j \gets 1$ \textbf{부터} $M_i$}
        \IF{$j > f_{S_i}$}
            \STATE $R'_j \gets \max\{0,\ 1 - ( j - f_{S_i} ) \times 0.1\}$ \COMMENT{요청 $j$가 가용량 $f_{S_i}$를 초과하면 부족한 GPU당 0.1 감소}
            \STATE $R_j \gets \max\{R_j,\ R'_j\}$
        \ELSE
            \STATE $R_j \gets 1$ \COMMENT{가용량 충분, 완전 적합}
        \ENDIF
    \ENDFOR
\ENDFOR
\STATE {}
\STATE $\tilde A_i \gets A_i - \min_k A_k,\ \forall i$ \COMMENT{상대 연령 재기준화}
\FOR{ $i \gets 0$ \textbf{부터} $l_{wq}-1$}
    \STATE $P_i \gets 1/\text{base}^{\,i}$
    \STATE $P_i^* \gets P_i + \alpha\, \tilde A_i \cdot R_{a_i}$
\ENDFOR
\STATE \textbf{반환} $P^*$
\end{algorithmic}
\end{algorithm}

\begin{algorithm}
\caption{SAFA 큐 재정렬}\label{alg:reorder}
\textbf{입력:} \\
\hspace*{4mm} 도착 순서의 대기 큐 $WQ=[\,t_0,\ldots,t_{l_{wq}-1}\,]$ \\
\hspace*{4mm} 알고리즘~\ref{alg:priority}의 조정 우선순위 $P^*$
\begin{flushleft}
\textbf{출력:} 갱신된 대기 큐
\end{flushleft}
\begin{algorithmic}[1]
\STATE $(P_{i_{\max}}^*, i_{\max}) \gets \max(P^*)$
\IF{$i_{\max} \neq 0$}
    \STATE $WQ \gets [\,t_{i_{\max}}\,] \,\Vert\, WQ[0{:}i_{\max}) \,\Vert\, WQ[i_{\max}{+}1{:}l_{wq})$ \COMMENT{$\Vert$: 연결; 슬라이스는 반개구간}
\ENDIF
\STATE \textbf{반환} $WQ$
\end{algorithmic}
\end{algorithm}

SAFA는 배치 직전에 실행되며 대기 큐와 현재 서버 풀의 할당 상태 외에는 읽지 않는다. 따라서 핵심 스케줄러를 건드리지 않는 배치 전 재정렬 계층이다. 단, 스케줄러 자체가 별도의 큐 스케줄링을 수행하거나 누적 메트릭 또는 맥락 정보에 의존하는 경우는 예외다. $P^*$는 작업당 한 번 계산되므로 한 패스의 복잡도는 $O(l_{wq})$이고, $R$ 계산은 서버 수에 선형이다.

\subsection{튜닝 없이 계수 $\alpha$ 설정}\label{sec:coeff}
이 절에서는 $\alpha$가 제어하는 절충을 정의하고, 워크로드가 바뀌면 최적값도 크게 움직여 고정값 하나로는 충분하지 않음을 보인 뒤, 큐 통계에서 값을 읽어오는 튜닝 없는 계수 $\alpha_{\text{eff}}$를 유도한다.

\subsubsection{절충 분석과 메트릭 정의}
$\alpha$가 지나치게 크면 막힌 작업을 너무 공격적으로 추월시켜 새로운 불공정을 만들 수 있으므로, 값을 선택하려면 무엇을 교환하는지 정량화해야 한다. 요청 가속기 수별로 누적 대기시간을 묶고 재정렬 전후의 변화를 누적 거리(accumulated distance, AD)로 정의한다. $C_k$를 $k$개 가속기를 요청한 작업 클래스, $WT_k=\sum_{j\in C_k} w_j$를 재정렬 전 총 대기시간, $WT'_k$를 재정렬 후 총 대기시간이라 하자. 가중치 $1/2^{M-k}$는 큰 요청 클래스에 더 큰 중요도를 부여한다.
\begin{equation}
AD = \sum_{k=1}^{M} \frac{(WT'_k - WT_k)}{2^{M-k}}
\label{eq:07}
\end{equation}

같은 가중치를 클래스별 최악 대기시간 변화에 적용하면 누적 최대 거리(accumulated maximum distance, AMD)가 된다.

\begin{equation}
AMD = \sum_{k=1}^{M} \frac{\max\limits_{j\in C_k} w'_j \;-\; \max\limits_{j\in C_k} w_j}{2^{M-k}}
\label{eq:08}
\vspace{2pt}
\end{equation}

여기서 $w_j$와 $w'_j$는 각각 재정렬 전후의 대기시간이다. AD는 대기시간 전체의 편향을, AMD는 각 클래스의 최악 대기시간을 포착하므로 둘을 동시에 최소화하면 평균적 공정성과 기아 억제를 함께 겨냥할 수 있다. 두 값은 makespan과 함께 식~\eqref{eq:09}에 들어가며, 각각 min--max 정규화되어 $\alpha$ 탐색용 목적함수를 구성한다. 가중치나 정규화가 바뀌면 최적 $\alpha$도 움직인다는 사실 자체가 고정 계수에 반대하는 추가 근거다.

$\alpha$는 0.01부터 1.99까지 0.01 간격으로 탐색하고, min--max 정규화된 makespan, AD, AMD 세 값의 합을 최소화하는 값을 선택한다. makespan은 트레이스의 마지막 작업이 끝나는 시각이다.

\begin{equation}
Obj_i^{\text{SAFA}} = makespan_i^{norm} + AD_i^{norm} + AMD_i^{norm}
\label{eq:09}
\vspace{2pt}
\end{equation}

\subsubsection{고정 계수의 한계}\label{coeff:limit}
고정 $\alpha$의 한계는 최적값이 워크로드에 따라 바뀐다는 점이다. 정규, 지수, 로그정규, 감마, 베타, Weibull, 균등, Poisson 등 여덟 분포 각각에서 1,000개 작업을 합성하고 식~\eqref{eq:09}를 최소화하는 분포별 최적값을 찾았다(표~\ref{fig:motivation}). 최적값은 감마 및 Poisson의 0.22부터 로그정규의 1.63까지 넓게 분포한다. 모든 워크로드에 공통으로 적합한 하나의 값이 없으므로, 어떤 고정 $\alpha$도 특정 워크로드에 맞춘 절충일 뿐이며 워크로드가 바뀌면 그리드 탐색을 다시 해야 한다.

\begin{table}[htb!]
\centering
\caption{분포별 최적 $\alpha$}
\label{fig:motivation}
\footnotesize
\setlength{\tabcolsep}{4pt}
\resizebox{\columnwidth}{!}{%
\begin{tabular}{@{}l*{8}{c}@{}}
\toprule
분포 & 정규 & 지수 & 로그정규 & 감마 & 베타 & Weibull & 균등 & Poisson \\
\midrule
최적 $\alpha$ & 0.53 & 0.44 & 1.63 & 0.22 & 0.26 & 0.55 & 0.50 & 0.22 \\
\bottomrule
\end{tabular}}
\end{table}

\subsubsection{튜닝 없는 계수의 유도} \label{alphaeff}
식~\eqref{eq:01}의 계수 $\alpha$는 우선순위 조정값으로 작업이 얼마나 쉽게 추월할지를 정한다. 우선순위는 선두의 1에서 0에 가깝게 감소하므로, 선두를 추월하려면 최대 1 정도의 차이를 메워야 한다. 조정값이 1보다 훨씬 작으면 추월이 일어나지 않고, 훨씬 크면 어떤 작업도 쉽게 추월한다. 고정 $\alpha$는 이 균형을 유지할 수 없다. 백로그가 쌓이면 $\tilde A_j$는 계속 커지고, $R_{a_j}$는 그 순간 서버 풀 상태에 따라 달라지기 때문이다.

따라서 두 스케일을 큐 자체에서 얻는다. 평균 상대 연령 $A_{\text{ref}} = \max(1,\text{mean}(\tilde A))$, $\tilde A_j = A_j - \min_k A_k$는 현재 큐에서 일반적인 작업이 갖는 연령이다. 최소값을 빼기 때문에 모든 작업의 연령이 같으면 $\tilde A_j=0$이고 평균도 0이 된다. 하한 1은 연령이 서로 비슷할 때 평균이 계수를 무한히 크게 만드는 것을 방지한다. 적합도 $R_{\min}=\min_j R_{a_j}$는 대기 중인 요청 가운데 가장 적합도가 낮은 작업의 값이며, 풀 혼잡이 심해질수록 작아진다.

계수는 두 스케일의 역수로 설정한다.

\begin{equation}
\alpha_{\text{eff}} = \frac{1}{A_{\text{ref}}\,R_{\min}}
\label{eq:auto-alpha}
\end{equation}

이 계수는 $R_{\min}>0$이어야 한다. 재생한 모든 트레이스에서 $M\le8$이고 $R_{\min}\ge0.2$이므로 항상 정의된다. 이를 조정식에 대입하면 두 비율로 분해된다.

\begin{equation}
\alpha_{\text{eff}}\,\tilde A_j\, R_{a_j} \;=\; \frac{\tilde A_j}{A_{\text{ref}}} \;\times\; \frac{R_{a_j}}{R_{\min}}
\label{eq:alphaeff-define}
\end{equation}

첫 번째 비율은 이 작업이 평균보다 얼마나 더 오래 기다렸는지를, 두 번째 비율은 가장 들어맞지 않는 요청보다 이 작업의 요청이 얼마나 더 잘 들어맞는지를 묻는다. 풀 혼잡이 심할수록 후자의 차이는 커진다. 둘 다 커야 작업이 승격되므로, 나이만 많다고 기아 상태에서 빠져나오지는 않는다. 평균적 연령이고 가장 적합도가 낮은 작업은 두 비율이 정확히 1이 되어 선두 우선순위 1과 같은 기준에 놓인다. 다른 모든 작업은 이 기준점에 상대적으로 위치한다. 따라서 큐가 아무리 깊고 오래 대기했더라도 조정값은 외부 상수가 아니라 선두 우선순위 1에 맞춰 측정되며, 이것이 튜닝할 값을 남기지 않는다.

연령에 따라 우선순위를 높이는 개념 자체는 새롭지 않다. 운영체제의 priority aging, Linux CFS의 vruntime 균등화, Tiresias\cite{gu2019tiresias}, Slurm의 multifactor age factor 모두 같은 아이디어를 쓴다. 차이는 가중치 설정 방식이다. 첫째, 조정값을 갱 자원 적합도 $R_{a_j}$와 결합해 별도 예외 규칙 없이 HOL 원인 작업의 우선순위를 낮춘다. 둘째, 연령을 도착 횟수로 세고 큐의 최소값에 대해 재기준화하므로 시간 단위가 계수에 들어오지 않는다. 셋째, $\alpha_{\text{eff}}=1/(A_{\text{ref}}R_{\min})$은 외부 상수 없이 큐 통계에서 직접 계산된다.

\section{실험 평가}\label{sec:eval}

\subsection{이산 사건 시뮬레이션을 통한 규모 검증}\label{eval:sim}
Python 이산 사건 시뮬레이터를 구축하고, 최대 요청량이 노드 경계인 8 GPU인 약 11만 1천 개 작업의 전체 Philly 트레이스\cite{jeon2019philly}를 256, 512, 1024 GPU 클러스터에서 재생했다. 시뮬레이션에는 단일 노드 B200$\times$8 Kubernetes 클러스터에서 측정한 표~\ref{tab:overhead}의 작업별 오버헤드를 반영했다. FIFO 기준선, 여섯 개 순서 결정 정책(SJF, LAS, Kueue, EASY, Themis, Lucid), 제~\ref{eval:robust:placement}절의 일곱 개 배치 정책을 모두 직접 구현했다. 저자가 코드를 공개한 Tiresias, Kueue, FGD, KAI, Lucid는 원 구현과 교차 확인했다.

\begin{table}[htb!]
\centering
\caption{작업별 생명주기 오버헤드}
\label{tab:overhead}
\small
\resizebox{\columnwidth}{!}{%
\begin{tabular}{@{}lcccc@{}}
\toprule
\textbf{단계} & 스케줄링 & 시작(단독/혼잡) & 종료 & \textbf{합계} \\
\midrule
\textbf{오버헤드 (s)} & 0.5 & 1.5 / 3.5 & 2.5 & \textbf{4.5~6.5} \\
\bottomrule
\end{tabular}}
\end{table}

\noindent 일곱 개 메트릭을 보고한다. 공정성은 하위 1\% 순서 공정성 점수로, 100은 도착 순서가 완전히 보존된 경우이고 0은 완전히 역전된 경우다. 추월 비율은 이후 도착한 작업의 과반에게 추월당한 작업의 비율이다. 중앙대기시간과 최대대기시간은 각각 중앙 및 최대 큐 대기시간(초)이며, K$=10^3$, M$=10^6$이다. 활용률은 평균 GPU 할당률이다. Gini는 대기시간 분포의 불평등도를 나타내며 0이면 모든 작업의 대기시간이 같다. 꼬리비율은 99백분위 대기시간을 중앙값으로 나눈 값이다. 중앙값이 0이면 0으로 보고한다.

순서 공정성은 작업별로 정의한다. $o_k$를 작업 $k$보다 늦게 도착했지만 더 먼저 시작한 작업의 수, $\ell_k$를 작업 $k$ 이후에 도착한 전체 작업 수라 하자.
\begin{equation}
F_k = 100\cdot\Bigl(1 - \frac{o_k}{\ell_k}\Bigr)
\label{eq:order-fair}
\end{equation}
$F_k=100$이면 순서가 보존된 것이고 $F_k=0$이면 완전히 역전된 것이다. $\ell_k=0$이면 $F_k=100$으로 둔다. 보고하는 공정성은 $F_k$의 1백분위, 즉 가장 심하게 추월당한 1\% 작업의 점수이므로 높을수록 좋다. 이 정의는 도착 순서 위반을 불공정으로 다루는 Sabin과 Sadayappan\cite{sabin2005unfairness}을 따른다.

Gini 계수는 한 실행의 $n$개 작업 전체에 대해 계산한다. $u_1\le\cdots\le u_n$을 대기시간의 오름차순이라 하면, $u_i$는 $i$번째로 짧은 대기시간이고 순위 $i$가 가중치가 된다.
\begin{equation}
G = \frac{2\sum_{i=1}^{n} i\,u_i}{n\sum_{i=1}^{n} u_i} - \frac{n+1}{n}
\label{eq:gini}
\end{equation}
$G=0$이면 모든 작업이 동일하게 기다린다는 뜻이며, $G$가 1에 가까워질수록 대기가 더 적은 수의 작업에 집중된다.

\subsection{기준선 구현과 충실도}\label{eval:impl}
비교 정책은 세 그룹으로 나뉜다. 차단형 FIFO는 모든 개선을 측정하는 기준이다. 여섯 개 순서 결정 정책(SJF, Themis, LAS, EASY, Kueue, Lucid)은 어떤 작업을 먼저 처리할지 결정하며 SAFA의 직접 비교 대상이다. 배치 축(most-allocated, compact, round-robin, MCTS, FGD, KAI의 binpack과 spread)은 선택된 작업을 어디에서 실행할지 결정하며 조합 가능성을 검증하는 데 사용한다.
\begin{itemize}
\item FIFO (차단형): 작업을 엄격한 도착 순서로 dispatch한다. 따라서 선두 작업이 슬롯을 얻지 못하면 뒤 작업 전체가 멈춘다. 이는 제~\ref{Related}절에서 정의한 차단형이다.
\item SJF: 더 짧은 실행시간의 작업을 우선하며, 실행시간은 트레이스에서 oracle 값으로 가져온다.
\item LAS (Tiresias): 누적 GPU 점유 시간이 가장 적은 작업을 우선한다. 비선점 엔진에서는 대기 작업의 누적 서비스량이 0으로 유지되므로 순서 키가 정보를 갖지 못하고 정책은 비차단형 FIFO로 환원된다. 여기서 LAS는 그 퇴화 형태의 Tiresias를 대표한다.
\item EASY: oracle 실행시간을 사용해 FIFO에 선두 예약 backfilling을 추가한다. 본 연구는 선두 작업의 예약 시작 시각 $T$ 이전에 완료되는 작업만 앞당기는 보수적 변형을 사용한다. 전체 변형은 한 구성에서 공정성을 높이지만 전체 패턴은 바뀌지 않는다.
\item Themis: 완료시간 공정성 메트릭 $r$이 가장 큰, 즉 가장 불리한 작업을 우선한다. 공개 코드가 없으므로 라운드 솔버 없이 $r$로 정렬하고 실행시간 정보에 20\% 무작위 오차를 주입했다. 이 근사는 큰 작업을 구제하지 못해 기아에 불리하게 편향될 수 있다.
\item Kueue: cohort borrowing으로 여유 자원을 빌리면서 VC quota를 우선하고 비차단 방식으로 동작한다. 여기서 낮은 공정성은 결함이 아니라 처리량 우선 설정의 비용이며, StrictFIFO에서는 공정성이 개선된다. Philly 트레이스에는 15개 가상 클러스터가 있어 quota 공정 공유가 활성화되고, Kueue는 LAS처럼 비차단형 FIFO로 환원되지 않는다.
\item most-allocated: 먼저 가속기 속도, 그다음 유휴 GPU 수가 적은 순으로 노드를 정렬해 작업을 집약하고 큰 갱을 위한 빈 노드를 남긴다.
\item compact: 먼저 가속기 속도, 그다음 노드 인덱스 순으로 정렬해 풀의 앞쪽부터 채운다.
\item round-robin: 유휴 GPU가 많은 노드를 먼저 정렬하고 동률인 노드 사이에서 시작점을 회전시켜 작업을 풀 전체에 분산한다.
\item MCTS: 현재 대기 큐에 대해 Monte Carlo Tree Search를 수행하고 rollout으로 후보 노드에 점수를 준다. 큐를 볼 수 없으면 most-allocated로 퇴화한다.
\item FGD: 큐는 FCFS로 유지하면서 단편화 증가량 $\Delta F$가 가장 작은 노드를 선택한다. 단편화 인지 배치의 효과를 분리하기 위해 전체 GPU 단위 갱에 맞게 특수화했다.
\item KAI: 미할당 GPU 수에 따라 노드를 순위화한다. binpack에서는 오름차순으로 정렬해 작업을 집약하고 큰 갱을 위한 빈 노드를 남기며, spread에서는 내림차순으로 정렬한다. 공개된 노드 점수는 이 개수에 단조이므로 개수만으로 순위화해도 원래 선택을 재현한다. KAI는 아직 들어맞지 않는 갱을 위해 자원을 파이프라이닝하기도 하지만, 본 엔진은 이를 차단형 admission으로 근사한다. 생략한 예약 보존 backfill은 KAI의 처리량을 높일 수 있지만 순서 공정성을 높이지는 않는다. KAI는 GPU 종류 선호를 적용하지 않으므로 most-allocated와 달리 더 빠른 가속기를 우선하지 않는다.
\item Lucid: GPU 가중 SJF로 작업을 정렬하고 sharing score에 따라 패킹해 비선점 collocation을 수행한다. 이 트레이스에는 저자들의 profiler를 사용할 수 없으므로, 본 구현은 합성 활용률, oracle 실행시간, 두 작업의 sharing score로 인덱싱한 0.85~1.00 범위의 collocation slowdown 표를 사용한다. 이 표가 제~\ref{baseline-rationale}절에서 언급한 합성 왜곡을 제한하는 단일 스칼라 계열이다.
\end{itemize}

\subsection{주요 결과: 효율성과 공정성 지형}\label{eval:main}
FIFO 기준선, 여섯 순서 결정 정책, SAFA를 256, 512, 1024 GPU의 세 클러스터 구성에서 평가했다. 표~\ref{tab:sim-single}과 그림~\ref{fig:cliff}에 결과를 제시한다. 표에는 SAFA가 두 번 나타난다. SAFA(고정 $\alpha$) 행은 세 구성 모두에서 하나의 계수 $\alpha=0.13889$를 유지해, 한 번 맞춘 값을 그대로 둔 채 클러스터가 바뀌었을 때의 결과를 보여준다. SAFA 행은 매 패스마다 큐 통계에서 다시 계산하는 튜닝 없는 계수 $\alpha_{\text{eff}}$를 사용한다.

\begin{table*}[htb!]
\centering
\caption{주요 결과: 순서 공정성(공정성 및 추월 비율), 큐 지연, 자원 활용률, 분포 공정성(Gini 및 꼬리비율). 중앙대기시간과 최대대기시간은 유효숫자 두 자리로 반올림했으며, 본문의 백분율은 반올림 전 값을 사용한다.}
\scriptsize
\setlength{\tabcolsep}{2.5pt}
\resizebox{\textwidth}{!}{%
\begin{tabular}{@{}l rrr rrr rrr rrr rrr rrr rrr@{}}
\toprule
 & \multicolumn{3}{c}{\textbf{공정성}} & \multicolumn{3}{c}{\textbf{추월 비율 (\%)}}& \multicolumn{3}{c}{\textbf{중앙대기시간 (s)}} & \multicolumn{3}{c}{\textbf{최대대기시간 (s)}} & \multicolumn{3}{c}{\textbf{활용률 (\%)}} & \multicolumn{3}{c}{\textbf{Gini}} & \multicolumn{3}{c}{\textbf{꼬리비율}}  \\
\cmidrule(lr){2-4}\cmidrule(lr){5-7}\cmidrule(lr){8-10}\cmidrule(lr){11-13}\cmidrule(lr){14-16}\cmidrule(lr){17-19}\cmidrule(lr){20-22}
\textbf{정책} & 256 & 512 & 1024 & 256 & 512 & 1024 & 256 & 512 & 1024 & 256 & 512 & 1024 & 256 & 512 & 1024 & 256 & 512 & 1024 & 256 & 512 & 1024 \\
\midrule
FIFO & 100.0 & 100.0 & 100.0 & 0.0 & 0.0 & 0.0 & 5.8M & 1.6M & 2 & 13.0M & 3.0M & 99K & 92 & 92 & 77 & 0.36 & 0.35 & 0.81 & 2.2 & 1.9 & 43,596 \\
SJF & 0.0 & 30.8 & 98.9 & 5.8 & 1.0 & 0.0 & 880 & 27 & 2 & 12.3M & 5.6M & 211K & 97 & 94 & 76 & 0.95 & 0.98 & 0.95 & 13,336 & 99,216 & 5,164 \\
LAS & 0.0 & 5.2 & 99.6 & 8.1 & 3.8 & 0.0 & 1.5M & 130K & 2 & 11.0M & 3.1M & 128K & 99 & 98 & 76 & 0.45 & 0.71 & 0.92 & 7.4 & 22.7 & 11,942 \\
Kueue & 0.0 & 0.0 & 99.5 & 7.2 & 3.8 & 0.0 & 19K & 2.1K & 2 & 12.2M & 4.9M & 224K & 98 & 95 & 77 & 0.92 & 0.94 & 0.94 & 605.3 & 1,909 & 7,971 \\
EASY & 0.0 & 0.0 & 99.5 & 4.4 & 3.4 & 0.0 & 5.0M & 1.3M & 2 & 12.3M & 2.8M & 86K & 92 & 93 & 77 & 0.40 & 0.41 & 0.85 & 2.4 & 2.0 & 37,687 \\
Themis & 0.0 & 7.9 & 99.2 & 6.4 & 1.7 & 0.0 & 3.5K & 378 & 2 & 12.0M & 4.4M & 162K & 97 & 94 & 76 & 0.94 & 0.97 & 0.94 & 3,082 & 7,665 & 5,438 \\
Lucid & 0.0 & 97.0 & 100.0 & 3.4 & 0.0 & 0.0 & 2.2K & 0 & 0 & 8.3M & 818K & 2 & 99 & 85 & 48 & 0.94 & 0.93 & 0.75 & 2,179 & 0.0 & 0.0 \\
SAFA (고정 $\alpha = 0.13889$) & 88.8 & 79.0 & 99.9 & 0.0 & 0.0 & 0.0 & 4.4M & 1.2M & 2 & 11.2M & 2.5M & 94K & 98 & 97 & 77 & 0.39 & 0.36 & 0.82 & 2.5 & 1.9 & 36,752 \\
SAFA & 93.5 & 92.7 & 100.0 & 0.0 & 0.0 & 0.0 & 5.0M & 1.5M & 2 & 11.9M & 2.7M & 102K & 96 & 95 & 77 & 0.38 & 0.34 & 0.80 & 2.4 & 1.8 & 45,416 \\
\bottomrule
\end{tabular}}
\label{tab:sim-single}
\end{table*}

\paragraph{공정성과 활용률}
그림~\ref{fig:tradeoff}는 모든 정책을 두 축에 동시에 배치한다. 기존 정책은 절충 관계를 따라 갈린다. FIFO는 공정성 100을 유지하지만 HOL 차단으로 GPU 활용률을 포기하는 반면, 다른 정책은 큰 작업을 굶기며 활용률을 높인다. SAFA는 어느 극단에도 속하지 않는다. 256 GPU에서 여섯 개 순서 결정 정책은 모두 공정성 0으로 떨어지지만 SAFA는 93.5를 유지하고, 512 GPU에서는 92.7을 유지한다(그림~\ref{fig:cliff}). 추월 비율은 SAFA와 FIFO가 0.0\%인 반면 여섯 순서 결정 정책은 3.4\%--8.1\%다(표~\ref{tab:sim-single}). FIFO는 분할 불가능한 갱 때문에 활용률 약 7퍼센트포인트를 포기해 92\%에 머무는 반면, 예약 없는 backfilling은 97\%--99\%에 도달한다. EASY는 선두 예약 때문에 backfill 여지가 작아 92\%라는 예외를 보인다. SAFA는 높은 순서 공정성을 유지하면서 활용률 격차의 대부분을 회복해 95\%--96\%에 도달한다.

\begin{figure}[htb!]
\centerline{\includegraphics[width=\columnwidth]{Pic/Fig_cliff.pdf}}
\caption{부하 증가와 공정성 급락.}
\label{fig:cliff}
\end{figure}

\paragraph{분포 공정성}
공정성은 순서 위반의 꼬리만 읽으므로 256 GPU에서 전체 대기시간 분포도 함께 살펴본다. 여섯 순서 결정 정책은 Gini 계수에서 두 그룹으로 나뉜다. SJF, Kueue, Themis, Lucid는 대기를 소수 작업에 집중시켜 0.92~0.95를 보이는 반면, LAS와 EASY는 각각 0.45와 0.40으로 대기를 고르게 퍼뜨린다. 꼬리비율에서도 같은 구분이 나타나 첫 그룹은 605~13,336, 두 번째 그룹은 2.4~7.4다. 그럼에도 여섯 정책 모두 공정성은 0이다. 대기를 고르게 분배하는 것과 작업을 도착 순서로 처리하는 것은 서로 다른 속성이며, 분포 메트릭만으로는 순서 위반을 볼 수 없다. SAFA는 Gini 0.38, 꼬리비율 2.4로 FIFO의 0.36과 2.2에 가깝고 공정성 93.5를 유지하므로, 순서 공정성을 소수의 매우 긴 대기와 맞바꾸지 않는다.

\begin{figure}[htb!]
\centerline{\includegraphics[width=\columnwidth]{Pic/Fig_tradeoff.pdf}}
\caption{GPU 활용률과 순서 공정성의 절충.}
\label{fig:tradeoff}
\end{figure}

\begin{table}[htb!]
\centering
\caption{구성별로 다시 맞춘 그리드 탐색 최적 고정 $\alpha$와 튜닝 없는 계수 비교}
\label{tab:ablation-grid}
\footnotesize
\resizebox{\columnwidth}{!}{%
\begin{tabular}{@{}ll@{\hspace{0.8em}}rr@{\hspace{0.8em}}rr@{}}
\toprule
 & & \multicolumn{2}{c}{\textbf{최적 고정 $\alpha$}} & \multicolumn{2}{c}{\textbf{SAFA}} \\
\cmidrule(lr){3-4}\cmidrule(lr){5-6}
 & \textbf{GPU} & 공정성\,($\alpha^\ast$) & 중앙대기시간 (s) & 공정성 & 중앙대기시간 (s) \\
\midrule
 & 256  & 88.8\,(0.14) & 4.4M & \textbf{93.5} & 5.0M \\
                     & 512  & 90.6\,(0.9)  & 1.3M & \textbf{92.7} & 1.5M \\
                     & 1024 & 100.0\,(0.01) & 2    & \textbf{100.0} & 2    \\
\bottomrule
\end{tabular}}
\end{table}

\paragraph{튜닝 없는 효과}
표~\ref{tab:ablation-grid}은 튜닝 없는 계수와 그리드 탐색으로 찾은 최적 고정 $\alpha$를 비교한다. Philly 트레이스의 각 구성에서 $\alpha$를 $0.01$, $0.02$, $0.05$, $0.1$, $0.13889$, $0.2$, $0.3$, $0.5$, $0.7$, $0.9$, $1.0$, $1.2$, $1.4$, $1.6$, $1.8$, $2.0$의 16개 값으로 탐색해 가장 높은 공정성을 내는 값을 선택했다. SAFA는 88.8 대비 93.5, 90.6 대비 92.7로 앞서고, 백로그가 거의 없어 모든 정책이 100에 가까운 1024-GPU에서는 둘 다 100.0으로 같다. 핵심은 점수 차이가 아니다. 탐색 비용은 표~\ref{tab:sim-single}에서 드러난다. 거기서 사용한 $\alpha=0.13889$는 256 GPU의 16개 값 중 최적으로 88.8을 내지만, 같은 값은 512 GPU에서 79.0에 그쳐 해당 구성에 다시 맞춘 계수의 90.6보다 낮다. SAFA는 그 구성에 대한 사전 정보 없이 92.7에 도달한다. 제~\ref{sec:eval}-\ref{eval:robust:trace}절에서는 한 번도 맞추지 않은 워크로드에서도 같은 현상을 보인다.

\subsection{독립 트레이스로의 일반화}\label{eval:robust:trace}
지금까지의 결과는 하나의 트레이스에서 얻었다. 따라서 Philly와 계보가 다른 두 공개 트레이스에서 평가를 반복했다. 첫째는 Helios\cite{hu2021helios} 데이터센터(SenseTime, SC'21)의 Venus 클러스터 트레이스다. seed-42로 50\%를 서브샘플링해 8-GPU 노드로 구성된 80-GPU 클러스터에서 재생하고, 클러스터보다 큰 요청은 제외해 모든 정책이 배치할 수 있는 62{,}069개 작업을 남겼다. 따라서 모든 정책은 동일한 작업 집합에서 비교된다. 패턴은 반복된다. 모든 순서 결정 기준선의 공정성은 0으로 떨어지고 추월 비율은 5.1\%~17.5\%인 반면, SAFA는 추월 비율 0.0\%에서 공정성 80.0을 유지하고 고정 $\alpha$는 58.0에 그친다. 둘째는 Alibaba PAI\cite{weng2022mlaas} 트레이스(NSDI'22)다. seed-42로 16.4\%를 서브샘플링해 8-GPU 노드의 256-GPU 클러스터에서 재생하고 같은 제외 규칙을 적용해 모든 정책이 배치할 수 있는 120{,}079개 작업을 남겼다. 다시 모든 순서 결정 기준선의 공정성이 0으로 떨어지고 추월 비율은 5.8\%~13.4\%인 반면, SAFA는 99.9를 유지하고 고정 $\alpha$는 94.7이다.

표~\ref{tab:trace3}은 세 트레이스를 모두 제시한다. 각 트레이스는 하나의 클러스터 크기에서 실행되며, Philly는 256 GPU, Helios는 80 GPU, Alibaba는 256 GPU다. 순서 결정 기준선은 3.4\%--17.5\%의 작업이 추월당하고 모두 공정성 0으로 떨어진다. SAFA는 이후 도착한 작업의 과반에게 추월당한 작업을 하나도 남기지 않으며 Philly 93.5, Helios 80.0, Alibaba 99.9의 공정성을 유지한다. 각 트레이스를 실행하면서 해당 큐에서 계수를 직접 읽었으므로 트레이스 간에 값을 넘기거나 다시 맞춘 것은 없다.

\begin{table}[htb!]
\centering
\caption{독립 대규모 트레이스에서의 일반화}
\label{tab:trace3}
\scriptsize
\setlength{\tabcolsep}{4pt}
\resizebox{\columnwidth}{!}{%
\begin{tabular}{@{}l@{\hspace{0.6em}}rr@{\hspace{0.9em}}rr@{\hspace{0.9em}}rr@{}}
\toprule
 & \multicolumn{2}{c}{\textbf{Philly}} & \multicolumn{2}{c}{\textbf{Helios}} & \multicolumn{2}{c}{\textbf{Alibaba}} \\
\cmidrule(lr){2-3}\cmidrule(lr){4-5}\cmidrule(lr){6-7}
\textbf{정책} & 추월 비율 (\%) & 공정성 & 추월 비율 (\%) & 공정성 & 추월 비율 (\%) & 공정성 \\
\midrule
FIFO          & 0.0   & 100  & 0.0  & 100  & 0.0  & 100 \\
SJF           & 5.8   & 0    & 8.8  & 0    & 6.0  & 0 \\
LAS/Kueue     & 7~8  & 0    & 17.5 & 0    & 10.6 & 0 \\
EASY          & 4.4   & 0    & 9.3  & 0    & 13.4 & 0 \\
Themis        & 6.4   & 0    & 8.5  & 0    & 6.6  & 0 \\
Lucid         & 3.4   & 0    & 5.1  & 0    & 5.8  & 0 \\
SAFA (고정 $\alpha = 0.13889$) & 0.0 & 88.8 & 0.0 & 58.0 & 0.0 & 94.7 \\
\textbf{SAFA} & \textbf{0.0} & \textbf{93.5} & \textbf{0.0} & \textbf{80.0} & \textbf{0.0} & \textbf{99.9} \\
\bottomrule
\end{tabular}}
\end{table}

\subsection{배치 정책과의 조합 가능성}\label{eval:robust:placement}
SAFA는 배치 전에 동작하므로 어떤 배치 정책 앞에도 둘 수 있다. 일곱 정책, 즉 most-allocated, compact, round-robin, MCTS, FGD, KAI binpack, KAI spread 각각에 대해 FIFO 아래에서 정책을 측정한 뒤 SAFA 아래에서 다시 측정했다(표~\ref{tab:placement}, 단일 256-GPU 구성, 주요 결과와 동일한 일곱 메트릭).

FIFO 아래에서 일곱 배치 정책의 성능은 동일하지 않다. round-robin과 KAI spread의 활용률은 각각 89.1\%, 89.3\%이고, 나머지 다섯 정책은 91.6\%~91.9\%다. 같은 두 정책의 중앙 대기시간도 6.1M\,s와 6.2M\,s로 가장 길며, 나머지는 5.8M\,s다.

SAFA는 일곱 정책 모두에서 활용률을 높이고 대기시간을 줄인다. 활용률은 94.7\%--96.5\%로 올라 3.7~5.8 퍼센트포인트 개선되고, 중앙 대기시간은 11\%--18\% 감소한다. 그 대가는 순서 공정성이다. 공정성은 FIFO의 100에서 91.8~93.7로 낮아지고, Gini는 0.35--0.36에서 0.38--0.39로, 꼬리비율은 2.23--2.24에서 2.32--2.42로 이동한다. 그러나 이는 표~\ref{tab:sim-single}의 모든 순서 결정 기준선이 기록한 0에 비하면 작은 변화다. 이득은 FIFO 기준선이 가장 낮은 곳에서 가장 크다. KAI spread에서는 활용률이 89.3\%에서 95.1\%로 올라 5.8퍼센트포인트 개선된다(그림~\ref{fig:placehol}).

\begin{figure}[htb!]
\centerline{\includegraphics[width=\columnwidth]{Pic/Fig_placement_hol.pdf}}
\caption{배치 정책별 HOL 회복 효과.}
\label{fig:placehol}
\end{figure}

\begin{table}[htb!]
\centering
\caption{SAFA와 배치 정책을 조합하기 전후의 메트릭(단일 256-GPU 구성)}
\label{tab:placement}
\resizebox{\columnwidth}{!}{%
\begin{tabular}{@{}ll rrrrrrr@{}}
\toprule
\textbf{배치} & & 중앙대기시간 (s) & 최대대기시간 (s) & 공정성 & 추월 비율 (\%) & 활용률 (\%) & Gini & 꼬리비율 \\
\midrule
\multirow{2}{*}{most-allocated} & FIFO & 5.8M & 13.0M & 100.0 & 0.0 & 91.9 & 0.36 & 2.23 \\
                                & SAFA & 5.0M & 11.9M & 93.5 & 0.0 & \textbf{96.0} & 0.38 & 2.36 \\
\multirow{2}{*}{compact}        & FIFO & 5.8M & 13.0M & 100.0 & 0.0 & 91.8 & 0.36 & 2.23 \\
                                & SAFA & 5.2M & 12.0M & 93.1 & 0.0 & \textbf{95.5} & 0.38 & 2.32 \\
\multirow{2}{*}{round-robin}    & FIFO & 6.1M & 13.8M & 100.0 & 0.0 & 89.1 & 0.36 & 2.24 \\
                                & SAFA & 5.1M & 12.4M & 91.8 & 0.0 & \textbf{94.7} & 0.39 & 2.42 \\
\multirow{2}{*}{MCTS}           & FIFO & 5.8M & 13.1M & 100.0 & 0.0 & 91.6 & 0.36 & 2.24 \\
                                & SAFA & 5.1M & 12.1M & 93.7 & 0.0 & \textbf{95.8} & 0.38 & 2.35 \\
\midrule
\multirow{2}{*}{FGD}            & FIFO & 5.8M & 13.1M & 100.0 & 0.0 & 91.9 & 0.36 & 2.24 \\
                                & SAFA & 5.0M & 12.0M & 93.5 & 0.0 & \textbf{96.5} & 0.38 & 2.37 \\
\multirow{2}{*}{KAI binpack}    & FIFO & 5.8M & 13.0M & 100.0 & 0.0 & 91.9 & 0.36 & 2.23 \\
                                & SAFA & 5.0M & 11.9M & 93.5 & 0.0 & \textbf{96.0} & 0.38 & 2.36 \\
\multirow{2}{*}{KAI spread}     & FIFO & 6.2M & 13.8M & 100.0 & 0.0 & 89.3 & 0.35 & 2.23 \\
                                & SAFA & 5.1M & 12.3M & 91.8 & 0.0 & \textbf{95.1} & 0.39 & 2.42 \\
\bottomrule
\end{tabular}}
\end{table}

\subsection{실제 클러스터 검증(B200)}\label{eval:b200}
시뮬레이터는 큐를 모델링하지만 Kubernetes admission 및 binding 경로는 모델링하지 않는다. 따라서 K8s v1.31을 실행하는 단일 노드 NVIDIA B200$\times$8 클러스터에서 교차 검증했다. 워크로드는 Philly에서 500개 작업을 샘플링해 시간 축만 압축하고 나머지는 바꾸지 않고 재생한 것으로, 원본의 bursty 제출 패턴을 보존한다. 피크 시점에는 500개 중 440개 작업이 동시에 대기한다. 각 작업은 실제 GPU를 점유하는 sleep stub로 실행하며 생성, gating, binding 경로를 실제로 거쳤다. 모든 정책과 반복을 합쳐 6회 $\times$ 500개 작업 $=$ 3{,}000개의 실제 Pod를 실행했다. 이 실험은 admission과 binding 경로를 검증하며 계산 및 통신은 제외한다.

큐 동작은 실제 클러스터에서도 유지됐다. SAFA는 두 실행에서 공정성 82.8과 82.4, 추월 비율 0.0\%와 0.2\%를 기록했다. 도착 순서를 구조적으로 보존하는 차단형 FIFO는 각각 84.2와 89.5였다. 선두의 대기 여부와 상관없이 후속 작업을 계속 dispatch하는 제어 실험은 공정성이 3.8까지 무너지고 추월 비율이 5.0\%가 됐다. 이 제어 실험과 앞 두 정책을 가르는 차이는 선두가 자리를 유지하는지 여부이므로, 이 측정이 입증하는 것은 SAFA가 실제 admission 경로에서 안전하게 실행될 수 있다는 점이지 연령 승격이 단순 FIFO보다 낫다는 점은 아니다.

단일 노드는 보여줄 수 있는 범위를 제한한다. 여기서는 연령 카운터가 도착 순서에 대해 단조 증가하므로 SAFA와 차단형 FIFO가 거의 같은 시퀀스를 만든다. 또한 노드 간 단편화가 없으므로 자원 적합도 변화도 작다. 따라서 이 설정은 시뮬레이션이 측정한 단편화 회복보다 순서 결정 동작을 검증한다.

\subsection{한계}\label{eval:threats}
위 결과는 실제로 측정하고 모델링한 범위에 의해 제한되며, 수치는 절대값보다는 경향으로 읽어야 한다. 실제 클러스터 검증은 단일 B200$\times$8 노드에서 SAFA, FIFO, 제어 실험만 burst replay한 결과다. 모든 순서 결정 정책의 실제 측정 비교와 다중 노드 단편화를 포함한 end-to-end 측정은 향후 과제로 남는다. 네트워크 또는 작업 간 간섭을 모델링하지 않았고, 작업 유형별 처리량 차이도 조사하지 않았다.

이 결론은 정책을 절대적으로 서열화하지 않는다. 공정성과 활용률의 절충이 어디에서 나타나며 부하가 높아질수록 얼마나 가파르게 심해지는지를 보여준다.

\FloatBarrier

\section{결론 및 향후 연구}\label{conclu}

본 논문은 GPU 서버 풀에서 HOL 차단이 유발하는 무기한 배제를 방지하기 위해 순서 결정 단계만 수정하는 배치 전 재정렬 계층 SAFA를 제시했다. 세 트레이스(Philly, Helios, Alibaba PAI)를 80~1024 GPU의 이산 사건 시뮬레이터에서 평가했고, FIFO 및 여섯 순서 결정 정책과 비교했으며, 일곱 배치 정책과 조합하고 NVIDIA B200$\times$8 Kubernetes 클러스터에서 교차 검증했다. 최근 스케줄링 정책이 실행시간 추정, 성능 프로파일 또는 솔버를 요구하는 것과 달리 SAFA는 스케줄링 시점에 이용 가능한 큐 위치, 자원 요구량, 상대 연령, 현재 서버 점유 상태만 읽으며 다양한 스케줄러 위에 계층으로 얹을 수 있다.

256 GPU에서 여섯 순서 결정 정책의 공정성이 모두 0으로 떨어지는 반면, SAFA는 일곱 배치 정책에서 91.8~93.7을 유지했다. 256 및 512 GPU에서 FIFO 대비 중앙 대기시간을 7\%--13\% 줄이고 활용률을 92\%에서 95\% 또는 96\%로 높여 공정성을 처리량 손실 없이 확보했다. 같은 패턴이 Helios와 Alibaba에서도 반복됐다. 경합이 약할 때는 양상이 다르며 Lucid 같은 collocation 정책이 더 빠르면서도 더 공정할 수 있다. 부하가 높아질수록 순서 결정 기준선이 순서 공정성을 잃고 크기 선택적 정책이 큰 작업을 굶기며, 여기서 SAFA는 아무것도 예측하지 않고 두 차원을 함께 유지한 유일한 정책이었다.

이 균형의 핵심은 계수 $\alpha$다. 고정값은 탐색으로 찾아야 하고, 최적값은 워크로드가 바뀌면 크게 이동한다. SAFA는 이를 큐의 상대 연령과 서버 풀에서 읽은 자원 적합도 $R$로부터 직접 계산하므로 새 그리드 탐색 없이 워크로드 변화에 적응한다.

향후 연구는 두 방향이다. 본 연구는 GPU gang 트레이스만 다루었으므로 Standard Workload Format의 고전적 병렬 작업 트레이스로 평가 범위를 확장하면 HPC 배치 스케줄링 문헌과 직접 연결할 수 있다. 더 중요한 것은 B200 검증이 다중 노드 단편화를 재현할 수 없는 단일 노드에서 수행됐다는 점이다. 따라서 여러 노드로 검증을 확장하는 것이 가장 우선적인 다음 단계다.

\FloatBarrier
\section*{감사의 글}
저자들은 본 연구와 관련해 어떠한 기관으로부터도 재정 지원을 받지 않았다.

생성형 AI 사용 공개. IEEE 권고에 따라, 저자들은 본 논문 준비 과정에서 생성형 AI 도구를 보조적으로만 사용했음을 밝힌다. 사용 범위는 문장 다듬기와 언어 교정, 그림 생성, 실험 재현 스크립트, 결과 정리 지원으로 제한했다. 연구 개념의 수립, 알고리즘 설계, 실험 수행과 해석은 저자들이 수행했다. 저자들은 AI 보조 콘텐츠를 모두 검토·수정했으며 출판 내용에 대한 전적인 책임을 진다.

\bibliographystyle{IEEEtran}
\bibliography{sn-bibliography}
%% if required, the content of .bbl file can be included here once bbl is generated
%%\input sn-article.bbl

% Author biographies --- to add a photo, replace IEEEbiographynophoto with
% \begin{IEEEbiography}[{\includegraphics[width=1in,height=1.25in,clip,keepaspectratio]{photo}}]{Name}.
% Fill in the XXXX/XX placeholders with actual degrees and years.
\begin{IEEEbiography}[{\includegraphics[width=1in,height=1.25in,clip,keepaspectratio]{Pic/bio_kyunam_cho.jpg}}]{Kyunam Cho}~는 서울 LG유플러스의 연구위원이자 고려대학교 컴퓨터공학 박사과정이다. 2000년부터 2021년까지 삼성전자에서 소프트웨어 개발자와 소프트웨어 아키텍트로 근무했고, 2021년부터 2022년까지 현대카드에서 소프트웨어 아키텍트로 일했다.

주요 연구 분야는 고성능 컴퓨팅(HPC)과 HPC 기반 소프트웨어 플랫폼이다. 삼성전자에서 수천 대 규모의 GPGPU 기반 머신러닝 시스템을 구축·운영했다. 현재 관심 분야는 LLMOps, 연합학습, 퍼블릭 클라우드 기반 플랫폼의 카오스 엔지니어링이다. 머신러닝 관련 저서를 집필했으며 \textit{Computer Physics Communications} 등에 논문을 발표했다.

\end{IEEEbiography}

\begin{IEEEbiography}[{\includegraphics[width=1in,height=1.25in,clip,keepaspectratio]{Pic/bio_younghwan_jin.jpg}}]{YoungHwan Jin}~은 연세대학교 컴퓨터과학과에서 학사 및 석사 학위를 받았다.
연구 관심 분야는 분산 시스템과 클라우드 네이티브 아키텍처이며, 특히 분산 GPU 추론을 위한 자원 할당과 엣지 환경의 확장 가능한 머신러닝 운영에 관심이 있다.

\end{IEEEbiography}

\begin{IEEEbiography}[{\includegraphics[width=1in,height=1.25in,clip,keepaspectratio]{Pic/bio_heonchang_yu.jpg}}]{Heonchang Yu}~는 고려대학교 컴퓨터학과 교수다. 고려대학교에서 1991년 석사, 1994년 박사 학위를 받았고 1998년부터 교수로 재직 중이다. 2019~2020년 미국 University of Texas at San Antonio 전기컴퓨터공학과 방문교수를 지냈으며, 고려대학교 Distributed and Cloud Computing Laboratory를 이끌고 있다.

연구 관심 분야는 클라우드 컴퓨팅, GPU 가상화를 포함한 가상화, 분산 시스템, 결함 허용 시스템, 실시간 처리를 지원하는 엣지-클라우드 기술이다. 2007년부터 한국정보과학회 편집위원으로 활동했다.
\end{IEEEbiography}

\end{document}
